\documentclass[b4paper,7pt,twocolumn,landscape]{jsarticle}
\usepackage{info_literacy_exam}
\usepackage{array,booktabs,colortbl}
\usepackage[table]{xcolor}

\begin{document}

\twocolumn[
\begin{center}
\large{\textbf{情報リテラシー 2026年度 前期中間試験 解答}}
\end{center}
\vspace{2mm}
]

\subsubsection*{解答一覧（各2点・全50問・100点満点）}

\vspace{2mm}

\begin{center}
\renewcommand{\arraystretch}{1.6}
\begin{tabular}{|c|c||c|c||c|c||c|c||c|c|}
\hline
\textbf{問} & \textbf{正解} & \textbf{問} & \textbf{正解} & \textbf{問} & \textbf{正解} & \textbf{問} & \textbf{正解} & \textbf{問} & \textbf{正解} \\
\hline
1  & 4 & 11 & 2 & 21 & 1 & 31 & 2 & 41 & 1 \\
\hline
2  & 2 & 12 & 2 & 22 & 2 & 32 & 3 & 42 & 4 \\
\hline
3  & 3 & 13 & 2 & 23 & 2 & 33 & 3 & 43 & 2 \\
\hline
4  & 3 & 14 & 2 & 24 & 2 & 34 & 2 & 44 & 4 \\
\hline
5  & 3 & 15 & 3 & 25 & 2 & 35 & 2 & 45 & 3 \\
\hline
6  & 2 & 16 & 2 & 26 & 2 & 36 & 3 & 46 & 1 \\
\hline
7  & 2 & 17 & 4 & 27 & 3 & 37 & 3 & 47 & 2 \\
\hline
8  & 2 & 18 & 3 & 28 & 2 & 38 & 3 & 48 & 3 \\
\hline
9  & 3 & 19 & 3 & 29 & 3 & 39 & 3 & 49 & 4 \\
\hline
10 & 2 & 20 & 3 & 30 & 3 & 40 & 2 & 50 & 3 \\
\hline
\end{tabular}
\end{center}

\vspace{4mm}

\subsubsection*{各問の正解内容}

{\small
\begin{enumerate}
  \item LINE（学内システムに挙げられていない）
  \item ログアウトする
  \item 多要素認証（2段階認証）
  \item 16文字
  \item AXIOLE（アクシオーレ）
  \item 学校のアカウントでログインする
  \item @outlook.jp
  \item コミュニケーション
  \item アンケート・テスト作成
  \item チャットは1対1やグループの会話、チャネル投稿はクラス全体への発信
  \item Forms
  \item 投稿は簡潔・丁寧にし、スレッド返信を活用する
  \item カメラ・マイクのオフを確認する
  \item 事実・観測結果をそのまま記録したもの
  \item 気温25℃
  \item コストをかけずに同じ情報を大量に複製できる
  \item 形式性
  \item 残存性
  \item 経済性（5視点に含まれない）
  \item 記録メディア
  \item 高速・世界規模・双方向性あり（インターネット）
  \item 問題の発見
  \item Check（評価）
  \item 繰り返すことで継続的に改善する
  \item 複数のキーワードをすべて含むページを検索
  \item 一次情報
  \item 批判しない・自由に発言・質より量
  \item カード作成→グループ化→図解化→文章化
  \item マインドマップ
  \item 著作物が創作された時点で自動的に発生する
  \item 著作者の死後70年
  \item 意匠権
  \item 私的使用のために著作物を複製する
  \item 引用部分が自分の文章より多い（正しい引用条件に含まれない）
  \item 生存する個人を特定できる情報
  \item 自分に関する情報を他人に知られず自己決定できる権利
  \item 写真のファイルサイズから個人が特定される（リスクに該当しない）
  \item OS（オペレーティングシステム）
  \item 個人Microsoftアカウント
  \item 4桁以上の数字で設定する
  \item 英数字とハイフンで設定する
  \item セキュリティの修正や機能の改善を適用する
  \item SSIDとユーザID・パスワード
  \item Word・Excel・PowerPoint
  \item portal.office.com
  \item 学校のMicrosoftアカウント
  \item クラウド上にファイルを保存するサービス
  \item チームでのコミュニケーションや共同作業
  \item 有線LANの方が安定した通信でインストールできるため
  \item 自己判断でアンインストールする（最も不適切）
\end{enumerate}
}

\end{document}
